\documentclass[UTF8,11pt,fontset=fandol]{ctexart}
\usepackage[a4paper,margin=22mm]{geometry}
\usepackage{amsmath,amssymb,booktabs,longtable,array}
\usepackage[colorlinks=true,urlcolor=blue,linkcolor=blue,citecolor=blue]{hyperref}
\usepackage{enumitem}
\setlist{nosep}
\setlength{\emergencystretch}{3em}
\title{18号赛题：距离寻址指令调度的算法研究与实测报告}
\author{本地工程研究记录}
\date{2026年10月11日；核心冻结05:10（北京时间）}
\begin{document}
\maketitle
\begin{abstract}
本项目针对基本块内的距离寻址指令重排与COPY接力，建立严格输入解析、独立符号验证、可配置整数解释器和固定种子评测体系。研究重点是减少过早产生的活跃值、避免全量扫描，并在预算内保留已验证的最佳候选。本报告区分官网要求、本地模型、合成输入成绩以及尚未证明的优化空间。冻结版本通过133回归和10000固定种子小图；最终336大图产生269合法输出，另列容量证书、预算与启发失败。严格窗口4终端组达到独立COPY下界，但收益不推广到一般DAG，也不宣称官方平均COPY目标已达成。
\end{abstract}

\section{任务、机器模型与限制}
官网页面为\url{https://www.chaspark.com/#/races/competitions/1326690549003280384}，本次核实时间为2026年10月10日。基本块由GET、SET以及一输入或二输入的运算组成；正文举出ADD和MULW，图中另有SLL但未给完整语义。输入字段为\texttt{block\_id}、\texttt{bit\_width}、\texttt{instructions}；相对引用\texttt{t\#d}指当前指令前$d$条的输出。对于$b\in\{2,3,4,5\}$，允许的最大距离为$W=2^b$，即2 bit允许距离4。超距需要COPY接力，重排后功能必须一致。输出核心字段为\texttt{block\_id}与\texttt{optimized\_instructions}；严格导出模式移除研究统计字段。

官网给出的Level 0要求正确性，Level 1/3给出平均COPY占比不超过10\%/2\%的目标；完整评分权重与占比分母尚未明确。项目同时记录COPY/输出指令和新增COPY/原指令两种分母，不能将自建数据结果当作官方成绩。官网三指令输入可以解析，但示意输出中的COPY使MULW的两个操作数都来自第二个GET；独立回归拒绝该错误输出，不能直接照抄。

当前仅支持明确配置的GET、SET、ADD、MULW、COPY；未知操作符和字段拒绝。GET/SET默认保留同寄存器RAW、WAR、WAW依赖，不同寄存器访问可以交换。验证器逐寄存器比较访问轨迹。另提供保守模式，强制全部GET/SET保持原序，用于历史对照。该额外限制会制造失败，不能当作原题无解。

本地整数模型显式配置位宽、符号和wrap/error溢出；MULW低位截断是假设而非官方语义。符号验证保留操作符及有序操作数，不采用结合律、交换律或重新计算，也不依赖猜测溢出行为。SHA-256符号签名是工程等价检查，存在密码哈希的理论碰撞限制；不是形式化证明系统。本文研究时保存所有原计算，每个原节点执行一次，只允许新增COPY，不做别名合并和外部spill。

\section{与本题实质相关的重要工作}
\subsection{距离寻址ISA的直接相关研究}
Sugita等在CC 2023提出距离寻址ISA的正确且完备代码生成方法\cite{sugita}。STRAIGHT以指令距离指定操作数，研究包括满足约束的生成技术及模拟器评估。这是本题最直接的文献关系。其编译流程还涉及spill、SSA消除与控制流，当前工程仅处理单个DAG基本块和COPY，因此不能继承论文整体的完备性保证。作者机构页面及DOI提供准确元数据，作者组原文链接也列于参考文献。

\subsection{表达式树与寄存器压力}
Sethi和Ullman研究受限寄存器机器上算术表达式代码生成\cite{sethi}。本项目借鉴先安排需求更大的子树这一思想，为DAG计算静态需求估计，再以依赖后序给GET和消费链确定优先级。共享子图、寄存器副作用和指令年龄窗口改变了优化问题；这些优先级是启发式，并不具有该树模型的最优性结论。

LLVM的\texttt{ScheduleDAGMILive}维护活跃区间并跟踪寄存器压力，通用调度器组合多种启发式\cite{llvm}。工程启示是把活跃状态维护作为调度过程的一部分，用就绪集合和增量更新代替每步重新扫描全部节点。传统寄存器容量与本题年龄窗口仍不同：一个值即使占用容量很少，也可能因过旧而不可访问。

\subsection{计算图、存储交换与组合优化}
Hong和Kung的红蓝pebble模型用DAG和有限快存储分析计算与I/O\cite{hong}。本题也需要控制物化中间值，但窗口限制取决于指令年龄，且当前模型没有无限慢存储的load/store，因此不能直接套用其I/O下界。它提示需要区分真实容量障碍、寿命过长和调度选择错误。

Unison以约束编程联合寄存器分配和指令调度，并在真实架构和应用上评价质量与编译时间\cite{unison}。本文有限状态搜索只优化COPY计数，且当前负载全部是合成基本块；不能把其跨架构效果作为本项目收益。值得借鉴的是显式目标、可行性证明边界、时间预算和候选保留机制。

\section{当前算法构思与工程实现}
\subsection{基线、搜索与独立验证}
冻结版本v0.2保存在\texttt{experiments/baseline\_v02/}，附逐文件SHA-256。其稳定基线保持原序；ready基线组合最后使用释放数、最旧操作数和关键路径。窗口即将过期时插入COPY；连续仅COPY无进展或超过COPY预算时终止，并将启发式失败与数学证明分开。

精确oracle默认至多10个原节点，状态为已完成节点集合与活跃原值的年龄，原指令转移代价0、COPY代价1；允许提前COPY。统一代价搜索完整耗尽或到达最优终点才产生模型内结论。证明仅覆盖保持原DAG边、原节点一次执行、所选寄存器依赖、仅新增COPY、无别名合并或重新计算的状态空间。大图用有界beam，裁剪后不能声称全局最优。默认30000状态、100 ms；时间检查在工作单元之间，解析、验证和I/O成本仍需统计，预算不是操作系统硬实时保证。

\subsection{第一批：增量活跃值与需求驱动顺序}
旧ready基线每步扫描全部remaining并筛选就绪节点，时间含二次项。新实现将仍有未来使用的已物化值放入按最后位置排序的OrderedDict，过期值查询为常数时间；就绪堆通过版本标记丢弃旧优先级条目。消费者数只在执行后更新；从两个消费者减到一个时定位唯一剩余消费者并更新相关就绪节点。数据结构预处理和状态空间存储为$O(n+m)$，堆操作为$O((n+m)\log n)$，输出和COPY治理另外与输出长度成正比。静态需求排序在任意大入度情况下含排序成本；本题数据操作数最多两个，但寄存器顺序边可能更多。

需求估计对不重复的子值需求降序排列，取\[r(v)=\max_i(r(u_i)+i),\quad i=0,1,\ldots\]作为静态优先级信息，并按sink的依赖后序生成排名。pressure策略首先最大化净释放，定义为\[\Delta(v)=\#\text{最后使用操作数}-\mathbf1(\text{当前结果还有未来使用}) .\]demand策略首先按后序排名；incremental策略组合释放数和路径高度。即将过期且只有一个就绪消费者时，优先执行该消费者，否则接力。所有将作为输出的候选都由独立验证器重新解析和重建符号，验证器不访问调度器内部状态。

第一批auto在大于128节点时移除旧基线，实测造成质量退化，因此该方案未作为最终结论。第二、三批用stable\_fast和ready\_fast严格复现旧稳定和ready选择；后者以最后使用数、最旧操作数和高度作为优先级，COPY使年龄键变差时在出堆时惰性修复，最后使用数改善时立即更新。480组拓扑/窗口/种子输入逐条输出与旧ready一致，另500个小图逐条复现stable。ready\_fast复杂度还包含惰性修复数$R$，不能将$O((n+m+R)\log(n+m+R))$省略成所有策略的线性复杂度；高fanout的COPY可能多次触发相关键修复。

候选达到零新增COPY时，已达到非负成本下界，可以立即停止该目标的搜索；这不证明运行时间、总原COPY或硬件执行周期最优。已有候选且节点数大于10时，设剩余预算为$T_r$、总预算为$T$，可选beam工作预算上限为\[\min\{T_r,\max\{5,\min\{25,0.2T\}\}\}\quad\text{毫秒}.\]无候选时仍允许使用剩余完整预算寻找可行方案。

第三批还改变搜索的同成本顺序：COPY数量仍是第一关键字，在相同成本下优先已完成更多节点的状态。它不改变统一代价搜索的最优成本结论，但改变了有界时间内可到达的状态。精确模式可用年龄支配：相同完成集合的活跃节点身份相同，成本不高且每个年龄不大的状态能够模拟较旧状态的合法后续；仅拒绝新的被支配入队，保留旧队列和父链，不在beam启用该规则。该性质依赖COPY-only、一次执行、相同DAG模型；Pareto桶比较可能带来额外开销，仍仅用于小图。200固定小图对不启用支配与进展优先的穷尽基准，最优成本及无解结果一致。

\section{实测体系、硬件与成本}
使用Windows 11、AMD Ryzen 9 8945HX（16物理核、32逻辑核）、31.3 GiB内存，CPython 3.12.14。运行依赖为Python标准库；无GPU、LLVM或第三方运行包，当前实验串行单进程。算法研发和独立验证无需目标ISA硬件；若需要评估机器周期、功耗或执行器行为，仍需官方模拟器或硬件接口，当前COPY计数不能替代这些测量。

首阶段10000固定种子小图覆盖wide、branch、highfanout和同寄存器hazard，2--5 bit各625个/家族，节点数约5--10。v0.2稳定、ready、auto分别验证通过9930、9992、9998个，所有输出均通过独立符号及三种整数配置差分。剩余两个种子4851、4995仅证明所选模型无解。auto相对两种基线最好值在503图减少COPY，并救回5图；数据整体运行约17.18秒。额外审计发现保守模式11个模型无解中9个在hazard模型得到合法调度，说明副作用约束必须精确。

第一批新增4组回归，包含1000固定种子小图乘3种新策略、8种拓扑乘4个窗口以及输入COPY/重复操作数、零COPY终止；连同原回归共39项，实测11.133秒通过。后续增至48项回归，独立上下文验证、错误提升候选拒绝等检查全部通过（27.283秒）。大图初始语料包含平衡树、偏斜树、并行链、菱形梯、稀疏共享、蝶形网络、多寄存器hazard和宽fanout八类；目标规模32、128、512、2048，每类每窗口3个种子，共384个输入。实际节点数随结构略有差异，逐条报告保存实际节点数而非只写目标值。规则树和网络的种子主要改变ADD/MULW标签，稀疏共享、hazard和fanout家族的种子也改变边或寄存器轨迹；样本之间存在结构相关性，不具有独立真实应用抽样的统计含义。

对照驱动\texttt{scripts/research\_benchmark.py}交替执行策略顺序，保存完整输入JSONL、每输入SHA-256、整个语料哈希、失败类别、逐条时间和CSV/JSON汇总。每个成功输出均接受独立符号以及signed8、unsigned16、signed32数值差分，避免大图任意精度爆炸影响成本。候选自身验证和外部差分成本分别计时。每批预算250 ms，最多30000扩展状态。数据与冻结源码位于\texttt{experiments/r1,r2,r3}及\texttt{baseline\_v02,baseline\_r1,baseline\_r2}。

\begin{table}[htbp]\centering\small
\begin{tabular}{llrrrr}\toprule
批次 & 版本 & 成功/384 & p50(ms) & p95(ms) & p99(ms)\\\midrule
r1 & v0.2 &279&109.49&289.76&313.56\\
r1 & 增量组合 &291&81.04&281.80&318.92\\
r2 & r1冻结版 &291&94.55&282.63&321.97\\
r2 & stable\_fast组合 &304&41.46&266.87&339.35\\
r3 & v0.2 &278&121.44&266.25&307.34\\
r3 & r2冻结版 &304&41.95&263.64&282.09\\
r3 & ready\_fast及搜索改进 &304&39.58&259.06&280.48\\\bottomrule
\end{tabular}
\caption{所有输入的优化器墙钟分位数，包含失败；每批串行且交替策略次序。预算不是墙钟硬上限。}\end{table}

不同批次v0.2成功数会因预算、系统竞争变化，不能把跨批次成功率差当作算法差异；只能比较同批同输入。r1共同成功267图中，7改善、230持平、30退化，COPY总数增加3915，揭示删除大图原序基线的负面效果。r2相对同批r1，共同成功291图的COPY从32736降到27498，27改善、263持平、1退化，新增13个可行输出。r3相对同批v0.2，共同成功278图COPY从18035降到16580（减少8.07\%），20改善、258持平、零退化，新增26个验证通过输出；共同成功的逐输入加速比中位数为3.05。所有成功均独立验证，无错误输出。

\begin{table}[htbp]\centering\small
\begin{tabular}{lrrrr}\toprule
r3拓扑 & 共同成功 & 旧COPY & 新COPY & 加速比中位数\\\midrule
平衡树 &24&501&495&4.24\\
蝶形 &18&30&0&2.70\\
多寄存器hazard &48&1305&0&3.42\\
菱形梯 &45&6300&6300&2.49\\
并行链 &48&0&0&2.58\\
偏斜树 &36&0&0&3.03\\
稀疏共享 &36&6220&6106&3.06\\
宽fanout &23&3679&3679&3.37\\\bottomrule
\end{tabular}\caption{同批同输入且两版均成功的质量比较；不把新增成功的大图混入COPY减少量。}\end{table}

r3仍有80输入没有可行输出：57为预算耗尽、23为调度失败，均不宣称原题无解。新版本整体成功集合的等权COPY/输出平均8.71\%，旧版为8.20\%，这一变化包含新增成功图较难的分布效应；只在成功子集上求平均会遗漏失败，不应据此判断满足官方10\%目标。原语料SHA-256为\texttt{61bb2e4fa6128d0473575a770f455660ad4f0dd83db1ccd165b4b527cc24fce8}。r1/r2/r3总耗时分别132.22、155.92、133.79秒（含多策略独立验证）。本机还有其他研究进程，且前两批部分时间与回归重叠；数值是当前会话观测，不是隔离硬件上的严谨吞吐保证。

\subsection{留出复验、搜索消融与第四批}
另使用种子100001起的256输入，目标规模64、256、1024、4096，八类拓扑、四个窗口、每格两个种子，生成后未用于前三批参数选择。v0.2与r3分别139、199成功，新增60个输出；共同139图中3改善、136持平，无质量退化，COPY10259降至10252；加速比中位数3.60。全输入p50由252.95降至59.57 ms，但p95由271.61变为278.82 ms，尾部仍被预算失败支配。这支持规模性能改善，但留出COPY收益只有7条，不能推广8\%质量收益到所有分布。证据位于\texttt{experiments/holdout}，语料哈希为\texttt{34503fa3f32dfbb928184064981069b53b039ec862268ac39bf60834025484f2}。

对种子30001起的1000小图进行精确搜索消融，FIFO、同成本进展优先、进展优先加年龄支配三种配置全部独立验证并给出相同最优成本，耗时约4.1秒。记录如下：
\begin{table}[htbp]\centering\small
\begin{tabular}{lrrrr}\toprule
配置 & 总扩展状态 & p50(ms)&p95(ms)&p99(ms)\\\midrule
FIFO &91639&0.146&3.713&7.386\\
进展优先 &11121&0.088&0.213&0.761\\
进展优先+支配 &10710&0.104&0.258&1.073\\\bottomrule
\end{tabular}\caption{1000小图搜索消融；各图250 ms、30000状态，全部完成最优成本搜索。}\end{table}
进展优先减少87.9\%状态；年龄支配只进一步减少3.7\%，但桶比较使延迟上升，因此默认关闭，只保留显式研究开关。证据为\texttt{experiments/oracle/report.json}；不是把状态减少自动等同于总时间收益。

第四批对所有引用原本在窗口内的输入，直接标准化原顺序输出并独立验证，达到零新增COPY下界；输入已有COPY仍保留。稳定增量策略和ready增量策略跳过未使用的需求排名；pressure在ready增量策略前尝试。相同384输入下，r3和r4均304成功，全部共同成功图COPY相同32805；加速比中位数1.16，p50由39.43降至37.28 ms，p99由287.62降至284.89 ms，总实验134.99秒。该批主要是运行成本收益，不能声称COPY进一步改善。

\subsection{独立进程内存观测}
\texttt{scripts/research\_resources.py}以独立Python进程测量四种拓扑的128、512、2048、8192目标规模，固定4 bit。16个输出全部验证。Windows峰值工作集包含解释器和模块，tracemalloc峰值记录Python分配；后者不涵盖所有系统/native内存。8192规模结果为：
\begin{table}[htbp]\centering\small
\begin{tabular}{lrrrr}\toprule
拓扑 & 实际节点 & COPY & Python峰值(MiB)&工作集峰值(MiB)\\\midrule
并行链 &8190&0&14.14&58.25\\
偏斜树 &8192&0&13.11&55.24\\
稀疏共享 &8192&1884&18.46&74.74\\
多寄存器hazard &8192&0&9.03&48.71\\\bottomrule
\end{tabular}\caption{r4实现的独立进程观测，预算20秒用于测资源；不代表250 ms下同样成功。}\end{table}
开启tracemalloc后，这四个8192图的生成、优化、验证墙钟分别约2.43、2.28、5.65、1.77秒，不能作为未插桩吞吐成绩。记录包含各规模的实际节点数，证据为\texttt{experiments/resources.json}。普通CPU和百MiB以内内存足以开展本批大图研究，仍不能由这16例给出所有图的峰值保证。

\subsection{第五、六批：共享输入亲和性}
新增多输出共享根值与扇出菱形两类拓扑，规模64、256、1024、4096，各窗口两个种子，十家族共320输入。affinity策略在净释放相同时先安排就绪计算，再提高直接共同操作数已经物化的待产生输入优先级。计数随原节点物化增量更新，避免每次扫描所有消费者。这能优先完成共享同一个GET值的消费链，推迟其他共享根的物化。

第五批先单独测试该策略：241共同成功图有26改善、171持平、44退化，另外13个原可行图在该单策略失败，因此不作为全局替代。其主要收益在多输出共享根值：24共同成功图COPY12130降至2528，另8个图恢复可行；其他图上常无效甚至较差。第六批按结构而非生成器名称选择：至少两个GET各有不少于四个直接消费者，且多个SET输出时，将affinity作为组合候选，仍保留其他基线和最佳有效候选。

在种子300001起的320输入对照中，r4/组合版254/262成功，共同254图COPY63405降至53803（减少15.1\%），24改善、230持平、无质量退化，新增8个输出；整体中位延迟约54 ms、同输入加速中位约1.00，说明这是质量收益而非普遍时间收益。记录位于\texttt{experiments/r5-affinity,r6-portfolio}。多输出共享的36个新种子、不同链长留出共同图COPY36810降至7740；全96定向留出（含8192规模，预算500 ms）由84成功提升至96成功。但是一个8192节点fanout图因预算只保留较差的原序候选，COPY262变3393，种子400095，已保留完整哈希和记录。该负面结果不能用总成功数掩盖；组合受有限预算影响，不能保证所有输入质量单调改善。

进一步扩大oracle到48个12/16节点拓扑，各配置1000 ms、100000状态。FIFO42个完成、6个预算耗尽，进展优先及可选支配各48完成；共同完成的最优成本均相同。默认规模限制仍为10，扩展实验显式允许16，不能把有界实验外推为任意大图的完备求解能力。证据为\texttt{experiments/oracle-hard}。

\subsection{第七批：独立验证复用与候选成本界}
大图重复重建原输入符号及验证不改善候选会消耗预算。新增独立VerificationContext，仅从原始JSON重新解析并建立原表达式快照；所有可能替换最佳结果的输出仍独立解析、检查窗口、比较符号和寄存器轨迹。不访问调度器的positions、remaining、heap或done。重复原图检查只做一次；COPY不优于已验证结果的候选明确记为未验证弃选，不能把这些弃选列为正确输出。搜索无改进时保留已验证原候选，并附搜索状态；外部研究驱动仍再次以单独的verify入口检查最终输出，并做三种整数差分。

320输入上，相比同批r6，260共同成功中259持平、1改善，COPY59359降至57088，加速比中位数1.10，p50由70.70降至45.32 ms。两版都261成功，但有一例各自独有成功，说明预算波动并未消失。原图快照改为首个可行候选产生后延迟建立，避免可行性探索之前提前消耗验证预算。

相同96个8192以内定向留出复验中，r6和验证复用版都96成功，共同图COPY40013降至36882，减少3131，正好恢复前述种子400095的较好压力候选；没有新增质量退化。全输入p50由134.99降至88.10 ms，p95由588.12降至537.60 ms，p99由662.82降至594.50 ms。依然明显超过500 ms名义预算，验证与工作单元不可抢占是实际限制，不能给硬实时承诺。记录分别为\texttt{experiments/r7-validation}和\texttt{affinity-holdout-r7}。该留出复验使用延迟建立之前的快照版，后续延迟版需单独验证，不能混写为同一个测量版本。

\subsection{结构随机留出与严格容量证书}
为避免只有运算标签变动的同构骨架，新增不规则多输出和跨组共享：种子决定GET根数、分支数、分支长度、根映射、原顺序交错和部分跨根边。种子500001起的96输入中，r4与第七批延迟上下文版各90成功，共同90图COPY83331降至71149，42改善、47持平、1退化；不规则多输出48图51958降至39778，但跨组42共同图仅改善2条，说明亲和策略不能处理所有共享关系。单独affinity更好，但可能被先执行的高COPY原序方案挤占预算，故后续组合将其放在结构符合时的首位，首次预算不超过总预算60\%（至少25 ms），保留回退策略。记录\texttt{experiments/r8-irregular}。

对纯无共享树/森林，年龄窗口可放松为$W$个可任意保留的寄存器槽位。在每个原节点只计算一次、不允许spill/别名合并/重算时，树需求递推给出该放松机器的必要槽数；若大于$W$，原年龄窗口模型必然不可行。仅在GET寄存器各不重复、SET寄存器各不重复且与GET无交集、没有COPY或未知运算、每个数据值至多一个消费者时应用此证书；共享DAG或寄存器别名均拒绝套用该递推。它不证明官方原题在允许其他变换时无解。回归以32节点平衡树的需求5对窗口4，以及内部窗口2的小树对完整oracle验证；公开输入位宽仍限定2--5，没有修改题面参数。

\subsection{一键完整里程碑验证}
\texttt{scripts/research.ps1 -Count 10000 -Samples 1 -Round milestone-v03}已实际运行成功：51回归18.521秒通过，10000小图9998输出全部独立验证，余下2仍是模型内无解。该小图测试的总时延p50/p95/p99为0.279/0.783/1.484 ms，历史v0.2为0.697/1.246/2.586 ms；这是跨时间观测，系统竞争不同，只作为性能趋势。两个COPY分母均与历史相同，故本批主要改善小图时间，未提升这一数据集的质量。

同一脚本的192个十二家族大图对照，v0.2与新版110/161成功；共同110图COPY9279降至7383，22改善、87持平、1图增加一条COPY，新增51可行输出。六个纯树由预算/调度失败转换为上述容量证书；其余仍有18预算和7调度失败。逐输入加速中位5.17，全输入p50由253.32降至38.81 ms，p99仍由286.55变为291.32 ms。证据\texttt{experiments/milestone-v03}包含完整小图CSV/JSON、最小诱导子图反例以及大图输入、逐输入报告和源码哈希。该阶段不是官方成绩，也不是所有输入质量单调改进的保证。

\subsection{第十批：静态图复用与优化余量量化}
组合各策略复用不变的依赖、消费者、高度和需求排名；每次运行独立创建remaining、pending、位置、活跃值、堆和亲和计数，避免策略之间污染。固定原序若峰值逻辑活跃值大于$W$，可跳过该顺序的基线；这只说明固定顺序不可行，不能转换为整个问题无解。新增顺序轨迹只用于研究成本，不被功能验证器读取。53回归16.187秒通过，包含策略交替多次运行的新鲜/复用结果逐条相同、无共享状态泄漏、错误图实例/顺序拒绝、COPY下界不超过实际值。

固定原指令顺序下，若原值生产到最后使用跨度为$L_v$，其自身$c_v$个接力至少满足$L_v+c_v\le(c_v+1)W$，因此
\[c_v\ge\max\left(0,\left\lceil\frac{L_v-W}{W-1}\right\rceil\right).\]
可把各值的必要自身接力相加。另取在$L$个原步骤内持续活跃的$k<W$个不同值；年龄总和每原步骤增加$k$，每次接力最多减少$W-k$，总和可变化范围不超过$k(W-k)$，故
\[C\ge\max\left(0,\left\lceil\frac{kL}{W-k}\right\rceil-k\right).\]
工程取最长跨度的若干前缀求这个下界，并与逐值下界取最大。它限制的是当前原节点顺序、不同值身份、一次计算、COPY-only模型；重排可能降低跨度，不能作为全局最优COPY证明。报告同时记录实际COPY减该下界的差值，作为固定顺序中最多可能改善量的宽松上界；不保证所有差值都可消除。

相同192输入、冻结r9与静态图复用版均161输出验证通过，COPY全部相同41603。优化器p50由45.21降至38.63 ms、p95由257.86降至254.63 ms、p99由277.29降至257.43 ms，逐输入加速中位1.06，总对照46.53秒。六例容量证书、18例预算、7例调度失败保持不变。161份核心字段输出实际写为JSONL，重新读取后以不同初始寄存器整数和三种固定宽度模型独立检查，全部通过，用时5.12秒；记录输出哈希、原输入哈希与源码清单，不依赖运行时调度状态。证据为\texttt{experiments/r10-prepared}。

当前顺序COPY下界对梯形菱形8334条解释8331，稀疏共享4997解释4948，宽扇出7070解释6941，跨组共享9841解释9700；多输出共享3040、扇出菱形3043和不规则多输出3078均达到该顺序下界。因此这些输入较大的COPY成本主要需要改变原节点顺序来降低，不能只更换接力时机。平衡树2200条仅有1638条被此松下界解释，剩余562条也不能直接称为可实现收益。

\subsection{第十一批：固定顺序精确接力审计}
实现独立的有界一致代价接力搜索：固定已完成前缀$i$时，活跃值身份由最后一次消费唯一决定，状态只保存这些值的年龄。若有$k_i$个活跃值，年龄互异且落在1到$W$，该前缀状态数至多排列数$P(W,k_i)$，全空间至多$\sum_i P(W,k_i)$。枚举下一条原指令与任何合法提前COPY；仅完整搜索或成本界排除所有更好状态时报告该顺序最优。窗口4可较小，但窗口8的组合空间增长显著，不得据此前缀分解称为多项式通用算法。

新增300个小图顺序，与原有完整mask/值身份oracle加原顺序依赖链的参考模型对照，所有最优COPY/不可行结论一致，并独立符号、整数检查；56项回归16.31秒通过。重新保存192输入中161成功输出及其原节点顺序侧文件，核心输出与顺序文件各有SHA-256，不能由表达式别名猜测身份。该运行成功数不变，但失败类别为19预算、6容量证书、6调度失败，显示预算边缘类别有波动。

161份输出中120份已由下界证书证明该顺序最优；29份有下界空隙但窗口大于8不搜索；12份窗口4/8且有空隙进行各500 ms、100000状态搜索，4完成最优性审计、8预算耗尽。仅一份64节点跨组共享由17降至16 COPY；64节点平衡树仍需11条，尽管松下界为8；4096节点窗口4跨组共享4078条经57064状态、452.13 ms确认该顺序最优。全审计含外部检查10.81秒。固定顺序搜索收益低且成本高，保留研究工具，未并入默认组合。证据为\texttt{experiments/r11-traces,r11-fixed-order}；不外推为跨全部重排的最优性。

\subsection{第十二至十四批：顺序探索的负例与新树收益}
第十二批针对稀疏共享、不规则多输出、跨组共享的96个新种子，尝试净压力与年龄结合以及年龄优先。前者81输出通过，旧auto为82；共同81图仅1改善、59持平、21退化，COPY28974变29196。单策略中位时间较低，但长尾更差且功能可行率下降；年龄优先更差。因此两者保留显式实验策略，不进入默认组合，不能用不同成功集合的总COPY制造收益。记录\texttt{experiments/r12-age}。

第十三批从上述COPY余量转向原节点顺序：Sethi式需求相同的兄弟子树，若实际工作量不同，先计算大子树可缩短其结果等待后一个子树的时间。工作量在无共享森林中按不同数据边递推为子树节点数；在共享DAG上只能作为截断启发式，默认组合不使用此规则。新增随机两两合并的非平衡树与带随机长度单输入依赖链的树两家族；后者每条ADD仍保留两个相同操作数，不删原运算。种子改变合并边或链长，非只有运算标签随机。

96个64--4096规模输入上，旧auto和大小需求策略共同78成功，COPY31357降至26362（15.9\%），67改善、11持平、无观察退化；其余18是森林容量证书。对同压力单策略比较为69改善、9持平，COPY31359降至26362，但逐输入加速中位0.98，说明收益主要是质量而非更快的单策略。第十四批冻结旧组合，仅在严格无共享森林且存在相同需求、不同工作量兄弟时，把大小需求候选放入组合，并保留旧候选；首次调度最多使用总预算60\%以留回退空间。相同96图组合COPY达到上述减少，但p50由35.59升至38.19 ms。

另64个新种子、128/512/2048/8192规模留出，旧新共同49成功，COPY37839降至33139（12.4\%），45改善、4持平；15剩余图仍为容量证书。组合p50由43.23升至48.43 ms、p99由332.25升至488.41 ms，清楚体现多一个候选的成本及预算/验证粒度限制，不宣称该组合普遍加速。留出总20.40秒，源码和输入哈希分别保留在\texttt{experiments/r14-size-portfolio,r14-tree-holdout}。在这些有限分布无退化并不构成所有图的质量单调保证。

\subsection{第十五、十六批：森林直接发射与窗口工作量顺序}
严格完整森林可以直接按后序发射并维护有序活跃值，不再经过一般就绪堆。1000个随机森林逐条输出与一般大小需求策略相同；同64图留出，49共同成功的全部33136 COPY保持相同，逐输入加速中位仅1.01，p50由45.68变45.41 ms。相比先前报告33139的差异来自有限预算运行，不归因于这个内核。说明移除堆在整个组合中的收益很小，证据\texttt{r15-forest-linear}。

进一步在两种子树顺序均满足剩余槽位时优先工作量大的子树，即使其需求较低。实际COPY发射仍检查年龄窗口。64图共同49成功，13改善、33持平、3退化，总COPY33136降至33002，仅减少0.4\%。该策略保持为\texttt{forest\_work}显式实验选项，未进入默认组合；单策略时间与多策略组合不可直接归因比较，记录\texttt{r16-window-work}。

\subsection{第二个完整里程碑与计算模板}
\texttt{milestone-v04}已实际一键执行：59回归16.337秒通过，10000小图9998个输出均独立符号及三种整数检查通过，2个模型内不可行。小图p50/p95/p99为0.284/0.802/1.433 ms。十四拓扑的224图中v0.2/新版122/187成功，共同122图COPY12728降至7731（39.3\%），33改善、88持平、1例多一条COPY，新增65可行输出。12例森林容量证书、18预算、7启发式失败单列；全大图对照74.14秒。其数据哈希为\texttt{264a5ab9\ldots2354f1d}，完整哈希与源码清单在JSON，不能把后来不同拓扑批次并为同一成绩。

R17加入四输出点积、三点一维stencil、方阵乘法tile三个计算模板。它们有独立整数金标准；是合成骨架，并非应用采集trace。此三模板的种子只改变标识，不改变结构，因此不把重复种子称为独立结构留出。48个目标64--4096节点、四窗口输入上，旧新23/32成功，共同23图COPY970降至616，7改善、14持平、2退化，新增9输出；其余15预算、1启发式失败。目标规模与实际规模不同：矩阵dim=2实际24节点，dim=4实际160节点。stencil全部16成功，窗口4分别需15/63/255/1023 COPY，更宽为0；点积窗口8/16/32全部12成功，窗口4四例预算耗尽。矩阵只有4个成功，不能由成功集合的均值宣称已覆盖矩阵应用。记录\texttt{r17-kernel-shapes}。

R21另生成48个所有原数据节点均在SET依赖闭包内的稀疏DAG；旧新27/36成功，仍有1预算、11启发式失败。这是生成不同输入分布，绝非从已提交失败图删节点。历史\texttt{sparse\_shared}允许未被SET观察的算术叶，额外压力不完全符合严格“SET叶”的理解，两种分布必须分列。当前生成器合计十八家族，所有实际节点数、结构种子、失败和哈希均保存。

\subsection{第十八至二十三批：失败的独立归因与严格证明边界}
对里程碑的18预算和7启发式失败，先把预算扩大到1500 ms、100000状态，25例没有找到新的物理可行见证，总26.54秒。这仍不能证明不可行。随后分别建立更严格可审计的必要障碍：8宽扇出图为强制算术脊柱的根值跨度容量界，必要槽数9大于窗口4或8；11配对网络中每对两条运算消费相同两种父身份、内部值恰被下一个对消费、终端有独立SET。首次SET前活跃数等于已产生GET数加半完成运算对数；首个终端对第一条运算需要其$R$个根且产生半对，至少需要$R+1$槽。8根/16根网络的必要9/17槽超过对应窗口。这些严格结构条件逐项筛选，未匹配的普通DAG不套用此式。

四个稀疏图通过依赖闭合子图证明：删除闭包外原指令及只为其服务的COPY只会缩短保留值年龄，所以一个闭包无解可给整个原图在同一身份模型内的必要障碍。R20搜索至多16节点的60候选闭包，获得7个证明，10个仍未知；其中6个由独立无年龄完成mask容量枚举复核，另一个14节点闭包在无年龄放松中可行，但冻结v0.2的无剪枝完整年龄搜索也判不可行。R22扩大至32节点、12个候选、各1000 ms，另外找到2个证明。候选搜索有界，闭包未必最小；不能把没找到障碍视为可行。

R23默认求解器只加入廉价脊柱与严格配对证书。224同输入、冻结r22/新版都187成功且50476 COPY相同；失败成为31模型证书、2预算、4启发式失败，总48.74秒，p50由35.89变34.19 ms但p99由256.50变269.39 ms，长尾未改善。之后默认加入纯SET祖先脊柱投影证书，识别1024节点跨组图至少需要6槽、窗口4。前述较昂贵闭包枚举没有进入默认求解器；离线证明不改变在线预算与失败记录。

\subsection{第二十四至二十六批：分支序列、全交错容量审计}
对算术脊柱互不共享、仅共享纯GET根、GET/SET寄存器唯一且互不相交的多输出图，实现串行分支容量DP。状态为已完成分支集合，尝试下一个完整分支，使用根首次/末次消费区间计算逻辑峰值，再实际生成每条原指令和合法COPY接力。最多16分支，状态至多$2^{16}$；时间仍受预处理、原节点规模、实际接力数量和外部验证影响，不能保证所有图都在该预算完成。

新种子1300001起96图、三家族、64--4096规模，单策略87成功、auto90成功；共同85图27改善、51持平、7退化，总COPY46598变46419，只减少179条。它救回2图而失去5图，p99达660.31 ms，故保留为\texttt{branch\_dp}显式选项。5个失去的图是未使用GET使覆盖筛选过严：纯唯一寄存器读可以保留原指令并提前执行，值不必活跃。R26冻结修复前版进行同输入对照，87升92成功，87共同图COPY49112全部相同，5图恢复；全对照15.98秒。没有删除原GET或算术指令，仍有4个串行策略失败，不能认为分支串行化覆盖所有DAG调度。

最后两例跨组图另由独立容量枚举审计：状态是各脊柱已执行前缀，允许所有分支交错；纯GET可推迟到首次消费，其间独立GET状态因此可以收缩而不增加逻辑容量。根的已产生/仍有消费者集合由前缀决定。完整穷尽才给出必要放松模型不可行，达到状态/时间预算则明确未知。80个小图与通用完成mask容量oracle的可行/不可行结论一致。256节点9分支图穷尽5321状态，179.70 ms；1024节点12分支图穷尽3923状态，201.30 ms。两者都判窗口4不足，证据\texttt{r25-branch-capacity}保存输入及源码哈希。

该专用全交错空间上界为各分支长度加两种终止状态的乘积，仍为指数；无年龄容量枚举可行也不代表COPY年龄窗口可行。这是离线工具，默认求解器没有启用这个较昂贵DP。至此原25个失败均有当前模型内障碍，\emph{不是25个官方原题全局无解}：所有证明保留不同原值身份，原计算只执行一次，不允许别名折叠、重新GET、重算、spill或代数变换。即使两个有序算术表达式完全相同，也仍按不同身份计容量；允许同值别名可能突破这些界，应作为新的模型单独研究。在线质量、在线成本、离线解释能力三者分开报告。该阶段77项回归实际28.869秒通过，数量随后续研究继续变化。

\subsection{第二十七至三十批：组合持出与验证成本收益}
另提供\texttt{branch\_portfolio}实验组合：结构符合时先尝试分支DP，其首次预算最多占整个组合40\%，并保留原affinity和全部基线、独立验证后的最佳候选。R27同96图共同85成功，27改善、58持平，无观察COPY退化，总30918降至30713；两版各87成功，但2个图各自独有成功，不能称为可行性单调改善。p50由46.69升至64.92 ms。R28用新种子1600001、128/512/2048/8192规模、500 ms预算持出96图，旧新91/90成功；共同90图17改善、72持平、1退化，总63021降至62847，仅174条；另外1成功图转为预算失败，逐输入加速中位0.93。收益小且增加成本，故该组合没有替换默认auto。成本预算和系统竞争会改变个别成功集合，历史负例保留。

R29按剖析推进验证成本优化。旧验证器为原输入及每个候选重复构建完整DAG，只用于格式检查；新独立验证器以一次原始JSON扫描检查字段、ID唯一性、操作数数量、相对引用范围及符号轨迹，不构建或读取调度DAG。GET初始签名按寄存器复用，读已有SET结果时不计算无用初始签名。500小图的完整符号状态与冻结旧验证器逐项相同，错误字段/ID/引用等变异的拒绝结果交叉一致；非标量opcode现在以明确格式错误终止。未知字段和运算符仍拒绝。

相同96图、冻结r28与新版都91成功，所有64173 COPY相同，逐输入整体加速中位1.07；整体p50反升60.52至78.66 ms，p99从581.23降至548.64 ms，因此不宣称普遍整体加速。为隔离验证成本，R30取这91份严格序列化输出，旧新验证各3次，交替执行顺序，无调度器调用。验证中位加速1.83，p50由23.83降至13.29 ms，p95由137.00降至62.07 ms，p99由201.00降至101.35 ms，总18.31秒。证据保存输入及输出哈希、逐次时间、源码哈希；它是验证阶段收益，不能将1.83倍写成整个编译流程加速。

\subsection{第三十一、三十二批与第三个完整里程碑}
就绪堆每次求键都重新扫描至多两个不同操作数以计最后使用释放数。R31将它变成增量计数：共享静态图保存初始计数，某源从两个剩余消费者变为一个时，只更新那一个消费者。400小图乘九种优先级的输出、失败结论及原顺序与冻结旧实现一致。同288图旧新227/228成功，共同227图225持平、2改善（预算机会变化），COPY54737变54729；逐输入加速中位1.008，全输入p99却273.85变308.52 ms。该实现减少重复工作，但不夸大为显著整体性能收益。

R32把严格完整森林的直接发射扩展到pressure与demand，两种需求后序均已与一般就绪优先级验证；1000个五家族32--128节点图乘三策略逐条结果一致。共享值、未观察原计算、寄存器别名以及保守总内存顺序依赖均拒绝直接森林路径。新种子1800001、五家族、128--8192规模的160图，共同131成功、29容量证书，43162 COPY全部相同；逐输入加速中位1.006，p50由40.18变39.64 ms，说明仅移除就绪堆仍不是当前完整组合的主要瓶颈。记录\texttt{r31-freed-count,r32-direct-forest}。

\texttt{milestone-v05}在R30版本上已实际一键成功：82回归、10000小图9998输出全通过及2模型证书。小图总时延p50/p95/p99为0.235/2.358/6.569 ms，长尾相比前批明显波动，不能只列中位数。十八家族、四规模、四窗口、每格一个种子的288图，旧新155/231成功；共同155图COPY14678降至9322（36.5\%），41改善、111持平、3退化，新增76可行输出。退化为64节点扇出菱形11变12、stencil14变15、矩阵tile7变8，输入及源码哈希全部保留。32在线模型证书、16预算、9启发式失败；大图对照85.74秒，整体p50由253.54降至33.41 ms，p99由330.29降至258.21 ms。它扩展了评测分布，不与十四家族224图直接合并计算提升率。

\subsection{第三十三批：修复近窗口容量的人工COPY剪断}
新增第十九类\texttt{near\_window\_spine}：$W-2$个共享根被一条长算术脊柱反复消费，种子改变根消费循环与运算标签。逻辑活跃数可达到$W-1$，在32槽机器上合法，但需要大量接力。原默认8倍原节点数的COPY保护是人为限制，会提前拒绝可行输入。由于连续$W$次COPY而没有原指令进展时已经终止，一个成功轨迹每个原节点前最多插入$W-1$次COPY，故默认改为$(W-1)n$保护；显式传入旧因子8仍可重现原失败。仍由时间预算限制实际工作，不是无界增加输出。

32个新输入、四窗口、128/512/2048/8192规模、两种子、2000 ms预算，旧新26/32成功；共同26图237872 COPY完全相同，新救回6个窗口32输入。512原节点需6437/6422 COPY，2048需29386/29393，8192需121628/121629，优化分别约1.61/1.43秒；相同严格序列化输出重新读取并以三整数配置再次检查，32份全通过，用4.29秒。全对照34.44秒，p99由2080.27降至1612.44 ms。较小128图两版均成功，其COPY也少于8倍界。

新成功集合的平均COPY/输出75.07\%，新增COPY/原指令543.29\%；这是高压力正确性覆盖而非达到官方占比目标。成功集合不同，不能拿全部COPY237872变552767声称质量退化或改善。后续若坚持身份模型，可用寿命/容量下界分析不可避免的高复制；若放宽同值别名等变换，须另建验证及证明范围。证据\texttt{r33-near-window}包含全部输入、CSV/JSON、输出及重新读取检查。

\subsection{第三十四至三十九批：森林预算反例、受限快路径与别名边界}
R34完整纯森林先执行，同160图仍131成功、29模型证书、43162 COPY相同，p50仅39.93变39.78 ms，总40.10秒。新增第二十类单源串行链森林：每个非SET原值只有一个不同消费者，多个独立GET在输入先集中出现，输出允许逐链执行。R35新48图（128--8192节点）全0 COPY，中位加速1.34，p95由295.76降至191.26 ms，用12.55秒。扩大新种子、256--16384节点的R36持出48图，却旧42成功、新40成功，共同40图3852 COPY变0；2个旧成功转预算失败，另6个两版均预算失败，总30.52秒。需求排序、通用图预处理与验证竞争成本不能被零COPY收益掩盖。一般森林优先因此移到显式\texttt{forest\_portfolio}，默认恢复r33顺序。

针对瓶颈另实现严格\emph{单源完整纯链森林}直接发射：沿每个SET回溯唯一不同输入，反向连续发射，重复二操作数都引用前一输出。仅在完整纯森林证书需求为1、无顺序hazard时启用；每条原节点恰好一次。它绕过通用图预处理和需求排序，时间$O(n)$、输出$O(n)$、0 COPY达到保留原节点模型的非负下界。100新小链独立检查全部原节点，保守总GET/SET順序模式禁止此路径。R38复用R36输入：本轮旧47/新48成功，共同47图17926 COPY变0，44持平、3改善，中位加速2.01，p99由712.92降至367.21 ms，总25.65秒。旧版同输入本轮47而前轮42成功说明系统竞争与预算敏感，不能跨轮把成功变化全归因于算法。R39新种子2400001起、512/2048/8192/16384节点、每格两种子的32图持出，旧30/新32成功，共同30图均0 COPY，中位加速2.06，p50由132.29变31.60 ms，p99由622.91变440.09 ms，总21.00秒。证据\texttt{r38-chain-direct,r39-chain-holdout}，策略不依据生成家族名或种子选择。

R37显式\texttt{alias\_pressure}探索身份模型边界。每个原GET和SET各为独立原子，重复同寄存器GET也不合并；仅COPY关系以及相同opcode、相同\emph{有序}算术子表达式共用最近输出。保留全部原指令、原DAG依赖与同寄存器RAW/WAR/WAW，不增加GET、不删除、不交换操作数、不作结合律变换。消费者仍等其原生产者执行，再可引用已物化同表达式值。它假设本地ADD/MULW纯运算，官网未明确授权替换操作数；是否用于正式模型待确认。GET符号或地址相同不足以推导额外读取允许，实验不依赖这种许可。

新模型旁路所有身份容量证书和固定顺序身份COPY下界。独立验证器仍仅读原JSON，以有序符号轨迹、计算多重集、GET/SET轨迹和三整数配置检查。扫描ready节点的原型最坏$O(n^2+nC)$，未取代堆。新增第二十一类同有序表达式配对网络，连同蝴蝶/稀疏共72新输入，64/256/1024规模、500 ms预算，总20.19秒。旧auto40输出、实验41，但共同35图21 COPY退化、14同，786变3056，中位加速0.42；4个旧输出转启发失败、1个转预算失败，没有平均质量收益，未纳入默认。6份身份模型容量证书输入在别名模型下获得正确输出：60节点窗口4新增COPY24/12，248节点窗口8为353/373，1016节点窗口8为1458/1537。它们说明证书限界必须清楚，\emph{不是旧证书在原模型内错误，也不是官方授权别名的证据}。见\texttt{r37-ordered-aliases}的哈希、输出、负例和成本。

\texttt{milestone-v06}在一般森林优先阶段完成87回归（53.05秒，保存日志）、10000小图9998正确输出与2模型证书，以及20家族320大图。旧新173/262成功，共同173图COPY21092变15976（24.3\%），41改善、130同、2退化，新增89输出；32在线证书、17预算、9启发失败。大图对照103.90秒，整体p50由254.03降至34.34 ms，p99由363.87降至280.28 ms。近窗口类产生极多合法接力，全部成功图COPY总数不直接比较。该里程碑随后依R36调整，历史源码哈希不覆盖新版；当前链快路径和显式别名另起完整里程碑复测，不能混写成该87回归版本。

\subsection{第四个完整里程碑与交付边界}
\texttt{milestone-v07}已一键完成92项回归（35.73秒）、10000固定种子小图（9.42秒）与21家族336大图（99.52秒）。小图9998输出全部符号和三整数配置独立差分通过、2身份模型证书；平均COPY/输出0.5386\%、新增COPY/原指令0.6138\%，只是本地小图分布。小图总时延p50/p95/p99为0.231/0.700/1.361 ms。大图旧新181/272输出，共同181图COPY19146降至16061（16.1\%），41改善、137同、3退化，新增91输出；39在线模型证书、16预算、9启发失败。全输入p50由253.24变32.69 ms，p99由347.37变257.49 ms，逐输入共同加速中位6.66。新成功集合平均COPY/输出14.02\%、新增COPY/原指令42.21\%，不能宣称官方10\%目标已经达到。

该336图输入哈希\texttt{7252b7ccaa...85d31b1c11}、逐图完整SHA256、严格输出和源码哈希保存于其目录。较大复制总量受近窗口压力样例及成功集合变化影响，质量百分比只在共同输出计算。R38/R39严格输出另外从文件重新读取，不调用调度器，48/32份均通过，5.91/4.96秒；R37实验41份也全部通过，用0.29秒。默认持续沿用身份模型证书，显式别名没有默认加分或掩盖在线失败。

最终审核还把一般森林候选的60\%首次预算限额也限定到显式森林组合，默认除单源链快路径外恢复r33候选/预算规则；这项变更与就绪扫描预算检查另保存源码。通用预算是协作式检查，解析与外部整数/符号验证可能超出调度预算，不能将500 ms参数称为端到端硬实时保证。高接力输出、较大共享图和真实应用分布仍是下一阶段瓶颈。

\subsection{第四十至四十二批：一般图复验与高接力分配成本}
R40用21家族、128/512/2048规模、各一新种子，252图与冻结r33对照，两版全部状态一致：203输出、30身份模型证书、11预算、8启发失败；98334 COPY全部相同，逐输入整体中位加速1.003，串行链子集2.00、parallel链子集1.34，长尾无普遍收益。它支持受限链路径之外的默认规则恢复，没有扩大成功集合的包装。

高接力输出至多用32种正距离字符串，却每条COPY都重新分配\texttt{t\#d}。R41每次调度建立最多33个不可变引用串，只复用字符串，不共享可变行/输入列表，序列化格式与语义相同。64新图、四家族、128--8192规模、2000 ms预算，旧新56输出、4证书、4启发失败；299632 COPY全部相同，中位加速1.01，p50由79.25降至73.93 ms但p99由1196.42升至1501.66 ms，无明显全流程时延收益，总31.95秒。

R42在8192节点窗口32、121624 COPY同输入，隔离进程交替旧新各三次，含生成/调度/独立符号检查的分配跟踪。所有6输出哈希相同。旧跟踪峰值为87.55/114.49/114.49 MiB，新为96.74/96.74/96.74 MiB；峰值中位减少15.5\%，进程工作集中位约315.98变264.88 MiB（16.2\%）。\emph{首对旧87.55低于新96.74，不能隐去}；该旧运行耗27.20秒，其余旧约12.8--12.9秒、新9.7--11.0秒，有限预算下尝试阶段与资源轨迹可能不同，首对成本异常未完整归因。资源跟踪时间不是普通吞吐测试；原31.95秒批量对照的长尾仍退化。证据\texttt{r42-token-resources.json}保留六次源码哈希、输入/输出哈希及原始观测，结论只限本工作负载分配中位收益。

\subsection{第四十三至四十八批：验证内存与数值差分成本}
验证器可在检查每个相对引用$d\leq W$且$d\leq$当前索引后，用$W$槽环形值历史代替整输出历史；连续ID已要求等于当前位置，故不必再存完整seen集合。原输入允许长距离，仍使用完整历史和ID集合。格式、合法引用、有序符号、计算多重集和寄存器轨迹逐原JSON重建，不读取任何调度状态。加入120小图的冻结全历史交叉检查、窗口边界、环绕、块首之前引用、当前/未来引用、未来/重复/布尔ID、未知PUT操作、SET寄存器篡改，以及50轮GET/ADD/SET副作用顺序拒绝测试。

R43取此前56严格输出，旧新各3次交替独立验证，中位加速0.999，基本持平；p99由192.03变188.38 ms，总9.96秒。R45隔离六次129821输出指令（8192原指令）的原始JSON验证，排除文件加载和调度，旧分配峰值每次9.15 MiB、新每次3.21 MiB（减少64.9\%）。跟踪时间旧约0.99--1.06秒、新1.14--1.38秒，仪器化成本反而增加，不能作加速成绩。环形表主要减少验证临时分配，整体求解峰值没有继续下降。R44另六次完整隔离进程观测的候选尝试序列相同：先验证stable再pressure、其余无改善丢弃、保留候选后的beam预算结束；COPY和输出哈希相同，旧均114.49 MiB、新均96.74 MiB，进程工作集旧约316 MiB、新约265 MiB。总95.59秒，补强字符串复用的资源证据，没有删除R42首对异常。

\texttt{milestone-v08}完成94回归41.12秒、10000小图9998正确与2证书，336大图旧184/新272输出；共同184图COPY21613降至16231（24.9\%），43改善、139同、2退化，新增88输出。39模型证书、16预算、9启发失败，99.26秒。整体p50由253.30降至31.27 ms，p99由359.90降至255.24 ms；所有272份严格输出重读独立符号及三整数配置验证通过，6.84秒。与v07共同图数量不同，不跨轮合并提升率。

进一步剖析高接力输出中的相对引用重复解析与整数传递：精确缓存规范\texttt{t\#1}至\texttt{t\#32}，其他串仍按严格词法拒绝或解析原输入长距离；不是跳过格式检查。整数初始寄存器与每个ADD/MULW结果已归一化，COPY/GET/SET仅传递已有值，归纳上无需重复取模/溢出检查。运算结果仍按配置归一化，unbounded乘法安全限界保留。500图乘五配置（有符号/无符号wrap、trap、受限unbounded）与冻结解释器的结果或错误完全相同；1000 COPY读写链与非法词法也覆盖。此为本地配置解释器成本优化，不推断官方机器字长。

R46同56严格输出的独立符号验证，旧新各三次，中位加速1.13，p50由14.67降至12.96 ms，p99由203.61降至166.78 ms，总11.50秒。R47同输出的三固定宽度整数模型，旧新各三次与原输入寄存器结果比对，中位加速1.82，p50由6.72变3.96 ms，p95由89.04变38.60 ms，p99由341.06变131.16 ms，总6.80秒。它们都是\emph{验证阶段}收益，不能写成求解全流程1.82倍；更广272份输出以及当前98回归、10000小图、336大图另起版本复测。

R48扩大到另一组21家族272严格输出，三整数配置、旧新各3次，全部与原输入一致；逐输出中位加速1.78，p50由1.439变0.754 ms，p95由20.179变10.814 ms，p99由64.332变33.082 ms，总10.13秒。该数据来自v08已成功输出，排除了未生成结果的64输入，验证成本结论不能当成整体正确率或优化质量提升。

\texttt{milestone-v09}完成98回归及10000小图9998正确/2证书；同336图旧183/新271输出，共同183图COPY21599变16231（24.9\%），42改善、138同、3退化，新增88输出，39证书、17预算、9启发失败，整体p50由253.25变30.87 ms、p99由362.18变258.23 ms。相较v08新版少一个4096原节点窗口32近窗口图，第一stable候选耗尽250 ms调度预算。该退化原输入和异常保留，不能用验证阶段加速宣称默认可行性单调增加；复制总数变化主要来自这类巨大输出失去，不能据此宣称减少了约6万COPY。

\subsection{第四十九至五十二批：小闭包在线证明与固定原序简化}
小矩阵窗口4的24原节点，用独立无年龄容量DP仅227状态完整证明当前身份模型不可行。较大图中可从原DAG识别二行二列的20节点依赖闭包：8个不同GET、8个两输入纯运算、4个合并节点。以原操作数有序结构匹配，不按矩阵名字选择；纯唯一不相交GET/SET与零内存顺序依赖筛选后，实际投影重新解析并完整容量DP。只对窗口4、至多4096组合与10 ms总识别/证明预算尝试；未知或预算直接没有证书。它首次进入默认求解器，成本计入整体解析/预检时间，和此前离线枚举不同。

R49四家族、64--4096规模、两新种子的128图，旧新均88正确、24194 COPY相同；8个矩阵失败转为实际在线模型证明（2原启发失败、6原预算），整体p50由56.01变36.86 ms，p99由540.09变530.23 ms，总52.82秒。没有新增可行输出，也没有把模型证明记成正确输出。

另识别三个共享三输入乘积前缀的27节点闭包：12个GET、9个纯产品节点、6个顺序合并节点；完整放松DP为1481状态。只窗口4、128组合、15 ms预算；保留原节点身份、操作数和每个GET，别名策略完全旁路。R51新128图、128/512/2048/8192规模，84正确输出和45507 COPY不变，15失败转在线证书，剩21预算、8启发失败，总59.35秒；第16个结构候选因识别与DP预算耗尽未证明，保留未知。小闭包状态上界仍$2^{27}$，这里实际穷尽1481状态才证明，不以10000状态预算当作完备性保证。

R50另在当时26个失败原输入上用2000 ms和100000搜索状态复测，总23.39秒，独立序列化见证救回前述近窗口4096图。其余按原错误分为4身份证书、8预算、13启发失败；历史脚本摘要\texttt{unresolved}只是没有可运行见证，不能忽略其中4证明。原报告保留，派生\texttt{classification-audit.json}记录源报告哈希和明确归因；以后脚本新增failure\_category并让任何独立校验阶段异常立即失败。

严格输入的每个依赖都指向更早位置，故稳定策略的最小ready总是下一个原节点。R52直接按固定原序发射，以相同到期/最后使用规则插COPY，取消就绪堆、优先级释放更新及额外状态；不是新搜索或新的COPY最优算法。800小图、寄存器hazard、21拓扑及原COPY/重复输入/旧8倍保护逐条与冻结实现一致。128新图、2000 ms预算，全部正确且577427 COPY相同，整体中位加速1.006、近窗口子集1.021，p95由583.85升至626.46 ms，p99由989.17变962.99 ms，总58.61秒；简化正确，但没有显著完整流程收益。

R53用集合视图包含判断消除每行字段检查的两个临时集合，非法/未知字段拒绝规则相同。五次50万次微测为1.43--1.53倍，而56实际严格输出、交替旧新三次完整符号验证只有中位1.017倍，p99由176.14升至187.14 ms，总9.31秒；不能从微测放大宣称流程收益。

最新\texttt{milestone-v10}106回归43.41秒、10000小图10.58秒（9998正确输出与2身份证书）、336大图97.54秒。大图旧183/新272正确，共同183图COPY21575变16219，41改善、139相同、3退化，89救回；47身份证书、9预算、8启发失败单列。整体p50由253.42变30.82 ms，p99由355.26变258.47 ms。相比v09证书增加不等于可运行解增加；源代码哈希和逐失败记录均保留。272严格JSON输出另行重新读取，不调用求解器，以不同寄存器初值和三个整数模型独立重放全部通过，总5.68秒。

\subsection{第五十四至五十九批：局部重排质量与一般图必要下界}
实现显式相邻交换\texttt{local\_swap}及半径至多$2W$、32位置的节点搬移\texttt{local\_relocate}。保持数据及同寄存器顺序依赖，只移动原节点；无需增加GET或替换算术。以跨度必要COPY下界的局部变化排序，最多512节点、128次试探、两轮、50 ms；先模拟逐条到期接力成本，只有严格少于已验证候选才发射，并由原始输入独立检查后晋升。时间或试探预算耗尽时保留已验证结果。空间$O(n+C)$（模拟$O(n)$），试探上界约$O(T(n+C))$，高扇出评分额外计邻接扫描；这不是最优顺序或最优接力算法。

R54六家族192新图，180共同正确、4证书、8启发失败；36改善、144同、零质量退化，COPY8257变8193，只减少0.78\%。p50由30.25升至38.70 ms，p99由52.08升至94.77 ms，总16.27秒。R55八家族、96/192/384/768新规模与种子的256图持出，218共同正确、29证书、1预算、8启发失败；29改善、189同，COPY13207变13161，只减少0.35\%，p50由30.60升至37.09 ms，p99由71.25升至89.95 ms，总20.03秒。两组保存严格输出独立重放180/218全通过；默认组合不启用。

R56跨位置搬移在相同192图上相较冻结相邻策略仅再改善4图、11 COPY，其余176持平，p50由39.62变44.75 ms，总20.38秒。R57同256图新结构持出，与自动策略相比33改善、185同，COPY13207变13154（0.40\%），正确数和失败分类不变；p50由30.25升至39.34 ms、p99由75.17升至92.00 ms，总20.73秒，218输出重放全通过。相邻邻域难带来可观收益；扩大移动仍不足以覆盖额外CPU代价，应保留实验并寻找子图级重排。

根据Bodlaender、Wolle、Koster的图收缩树宽下界研究\cite{minor}，另构造必要身份证书。任何活跃原值前沿至多$W$的拓扑顺序，令每个路径分解袋为“发射前活跃值加当前节点”，袋大小至多$W+1$，因此底层无向数据依赖图树宽至多$W$。图的任一minor最小度都下界树宽；故实际合法边收缩得到最小度大于$W$的图，即可反驳该身份模型。忽略寄存器顺序依赖只放宽问题。别名、重算、重复读取、spill仍不在证明范围；不得推广为官方全局无解。

R58朴素扫描收缩64图只作探索，较大图扫描约0.5--0.7秒；随后采用最小度惰性堆，并比较按邻居度及共同邻居数量的收缩策略。R59对v10剩余17失败原输入、每策略100 ms离线审计，总0.81秒，为9个给出独立有效minor：6个稀疏图、3个矩阵图。单次搜索约1.52--43.14 ms，独立检查约1.61--34.95 ms，二者分列。检查器不读启发式堆/邻接状态，而在原始JSON上以并查集维护收缩分支集合，再从原始边重建最终minor，验证每次收缩确有边并重新计算最小度；200小图与完整容量oracle交叉检查、篡改输入/窗口/度/收缩拒绝测试均通过。没有新的可运行见证；另外8失败仍未知。这批是离线证据，显式在线组合另行测量。

R60显式在线minor组合先在336固定图上尝试，272输出不变，仅1失败改证明。R61避免对每个提议重复符号检查：已解析DAG只能用于启发式提议，阳性结果仍从原始JSON独立重放；并去掉不必要的GET数量门槛，25 ms预算内新增6证书。仍272正确，共同图有1个近窗口COPY增加354，其余271相同，总195171变195525，中位加速0.990，总57.23秒，p99由259.53升至264.67 ms；时限减少影响候选到达顺序，不能声称预检没有质量代价。

R62六家族、96/384/1536/6144新规模与种子的192图持出，旧138/新137输出，16失败改独立证书、1输出失于预算。共同137图1改善、135同、1退化，COPY494929变494795，逐输入中位0.997；p50由73.09变61.01 ms、p99由570.25变542.86 ms，总77.21秒。成功集合变化与证明耗时分别记录；137结果严格重放另测。该组合保持显式，不加入默认。113完整回归47.93秒通过；前次一旧断言要求在线15 ms必须完成证明，在CPU竞争下返回合法未知而失败。修正为独立1000 ms完整投影证明与在线允许未知失败，正确输出检查没有放宽，失败日志说明保留。

\subsection{第六十三至六十五批：分支真实接力成本与脊柱优化余量}
局部搬移只带来不足1\%收益后，尝试更大粒度的完整纯共享GET分支顺序beam。限定至多16条互不共享算术的脊柱、唯一且互不相交GET/SET寄存器，保留全部原节点。状态记录已完成分支集合、每个仍活跃GET根的真实引用年龄；转移按实际窗口到期规则完整模拟一条分支的COPY成本，按成本与剩余活跃根排序，宽度16、至多4096状态。最终原序重新发射，并独立符号/数值验证。该有限beam既非一般图搜索，也非全局COPY最优；至少48图转移成本/发射一致性、原节点覆盖与三类拒绝路径测试通过。

R63的96新图，旧分支93正确、新87，初版给回退容量分支仅50 ms，丢失6旧有效输出。共同87图仅16改善、71同，COPY31480变31368，不能把丢失巨大输出导致的总COPY下降当作收益。R64修正为先让旧容量分支在共享预算内完成并独立验证，再以剩余时间改善；同96图93/93正确、3失败，15改善、78同、无质量退化，COPY36569变36459，仅0.30\%。p50由22.00升至165.92 ms，p99由263.36升至603.66 ms，逐输入中位速度0.145倍，总36.70秒；93严格输出换初值重放全部正确，1.94秒。没有加入默认；新的6144节点持出另测。

为判断高COPY图还能改善多少，单脊柱新增对\emph{所有拓扑重排}适用的必要COPY下界。筛选与强制脊柱容量证书相同，GET可自由延迟；算术链顺序由依赖固定。选择$k<W$个GET身份，其第一至最后强制链使用之间共同包含$L$个内部算术步骤。它们在整个区间都存活，引用年龄总和每个算术步骤增加$k$；任一COPY至多减少$W-k$，总和范围至多$k(W-k)$，故
\[
 C\ \geq\ \max\left(0,\left\lceil\frac{kL-k(W-k)}{W-k}\right\rceil\right).
\]
只枚举按使用跨度排序的根前缀，取最大有效下界；这不保证找到最强根集合，但每个证据成立。200个至多10原节点的脊柱与完整原指令/COPY年龄状态最优oracle交叉检查，均未超过真实最优；未知结构、寄存器副作用、多输出图拒绝。原COPY不在筛选范围，因此这里$C$同时等于新增和总COPY。

固定4096节点近窗口图，窗口32实60008 COPY、所有重排下界59820，因此在该身份模型中最多还能减少188条，即0.31\%；窗口4实4090、下界4088。这说明巨大接力数主要来自窗口容量，而非单纯启发式不够好；对这些合成极端图不能承诺2\%占比。它不代表官方工作负载或官方允许的全部重写模型，别名、重算、额外读取和spill仍不涵盖。新的64种子/规模批次继续测此余量。

R65新64图、128/512/2048/8192节点、四窗口、四种子，2000 ms预算下全部独立符号与三整数模型正确，总58.77秒。32个至少2048节点图，所有重排必要下界将可降低COPY比例限制在最多0.5987\%；8192节点窗口4/8/16/32四种子最大余量分别0.0244/0.0734/0.2083/0.4124\%。窗口32该规模的必需COPY/最终输出比例至少93.68\%，这是特定合成脊柱与身份模型的数学限制，不是官方平均分布成绩。完整输入、证书、源码哈希、计时及差分结果保存在\texttt{r65-spine-global-margin}。

R66整分支beam新96图持出，96/384/1536/6144规模、750 ms预算，旧86/新84正确，丢2预算案例，另6共同预算、4共同启发失败。共同84图20改善64同，无质量退化；COPY42645变42279（0.86\%），p50由29.57升至380.75 ms，逐输入中位0.147倍。总成本约72秒，84严格输出换初值重放全通过2.52秒。实现虽然保留\emph{已取得}的验证候选，但前置结构/验证工作仍可能使有限总预算无法取得旧候选，不能把回退保留说成每次都维持旧可行数。

最新\texttt{milestone-v11}118完整回归38.51秒、10000小图8.86秒、336大图97.40秒。小图9998正确加2身份证书；大图旧182/新272正确，共同182图COPY21601变16231，42改善137同3退化，90救回；47证书、9预算、8启发失败。新增实验均显式，不影响默认271/272等历史有限预算波动解释。默认当前成功集COPY总195171与v10一致；没有用实验不佳版本替代稳健默认。

\subsection{第六十七至七十批：带COPY余量界限的显式成本控制}
强制脊柱下界不仅说明质量余量，也允许明确的成本取舍。新增显式\texttt{bound\_portfolio}：只对至少512节点、筛选成立且正下界的脊柱，在已经独立验证的候选成本$c$、必要下界$b$满足$(c-b)/c\leq\varepsilon$时停止余下候选与beam。默认实验$\varepsilon=0.005$，表示\emph{当前COPY中至多还能减少0.5\%}；它不是COPY/输出比，也不是官方2\%目标，更不意味着已经精确最优。CLI支持\texttt{--copy-gap-tolerance}，0仅在精确达到下界时停止；不匹配图沿用普通组合。默认auto保持以质量为先，未启用该近似停止。

R67四新种子、512/2048/8192/16384节点、四窗口的64图，4000 ms共享预算，旧新64全部正确；36相同、28COPY退化，总2796368变2799745（增加0.121\%）。逐输入中位加速2.515，p50由348.69变182.57 ms，p95由2011.47变609.88 ms，p99由2462.46变789.87 ms，总87.97秒。R68三家族含非匹配图、768/3072/6144/12288新规模、两种子的96图持出，旧新均72正确、16身份证书、8启发失败；56同、16退化，COPY1244605变1245615（增加0.081\%），逐输入中位1.293，近窗口子集1.865，p50由178.69变140.03 ms，p99由1621.59变674.62 ms，总70.26秒。成本收益明确，但质量损失必须同时报告。

两批严格输出另换初值、三整数模型重放64/72全部正确，15.96/9.29秒；另由独立原始JSON检查器重建强制链、唯一寄存器、每个根的第一/最后使用区间与年龄和不等式，64/72符号输出正确，64/48证据通过，其中48/32结果已处于0.5\%余量范围，13.90/10.89秒。后者不调用调度器，区间选择和排序不重用启发式摘要；原COPY、额外消费者、错误区间/下界和未知结构拒绝。R67/R68上述求解耗时还不含这次离线证明复核，实际求解器在线加入复核后另作R70对照，避免隐藏检查成本。

\begin{table}[htbp]\centering\small
\begin{tabular}{lrrrr}\toprule
R69实验余量阈值 &证明早停数/18&总COPY&p50/ms&p95/ms\\\midrule
普通auto&0&979558&371.12&2035.26\\
0.1\%&7&979942&363.10&4214.48\\
0.2\%&10&979674&325.11&1193.65\\
0.5\%&14&980469&173.28&810.74\\
1.0\%&16&981156&135.82&3976.49\\\bottomrule
\end{tabular}\caption{18新图五配置成本曲线；全部输出独立正确。}\end{table}
R69为1024/4096/16384节点、窗口4/16/32、两种子，五配置共90调用，总95.03秒。所有阳性停止都由原JSON独立区间证书复核，非停止案例仅普通有界组合，没有额外近似最优保证。0.1\%和1.0\%的长尾异常保留：同机并行工作与有限预算候选阶段影响耗时，阈值变松并不保证这次小样本p95单调下降；不能选一个阈值宣称一般工作负载最优。0.5\%仍只是显式实验默认，不替代正式auto。

R70把原JSON证明复核放入真实求解过程，96图仍72正确、COPY1245615相同。相较不在线复核版本逐输入中位0.982，p99由907.89升至3401.65 ms，总84.05秒；最慢检查约191.58 ms，但多数差额发生在检查前的调度，无证书的稀疏图也由约907升至1878 ms，异常没有完整归因，不能全归因证书检查。该负面批次原样保留。

R71将证明检查改为完整的独立纯指令语法、ID、引用、寄存器与强制链/根区间验证；不再构建该数学证明不需要的第二个SHA符号图。候选输出原有独立符号与数值验证继续执行，未知字段/操作符、错误元数、无效寄存器、未来/当前/非规范引用、重复ID均拒绝。96同输入仍72正确、全部COPY相同，逐输入中位1.021、近窗口子集1.063，p99由727.28变732.99 ms，总65.98秒；这轮没有完整流程的大幅收益。R72在64原输入上交替旧新检查各三次，证据一致，证明检查阶段中位1.911倍，总11.95秒，不能写成整个优化器1.911倍。

R73新48图含不匹配随机树、1024/4096/16384规模和两新种子，含\emph{在线独立证书检查}的最终版本与普通auto对照：40/40正确、8身份证书，30COPY同、10退化，COPY1126318变1127281（增加0.086\%），整体逐输入中位1.194、匹配脊柱子集1.875、不匹配树0.987。p50由134.23变126.70 ms，p95由1240.42变763.67 ms，p99由1813.93变970.95 ms，总45.84秒。此轮才用于判断最终可执行成本与质量取舍；40份严格输出另换初值独立符号/三整数重放全部通过，9.01秒。之后121回归全部通过，37.83秒；默认auto算法没有因此改为近似停止。

R74另用新种子的16384原节点、窗口32，隔离进程普通auto与显式0.5\%停止各两次，生成、调度和独立符号检查均跟踪分配。四输出全部正确；普通244482 COPY，显式244821 COPY，增加339（0.139\%），证书证明显式结果最大可削减比例不超过0.205\%。普通分配峰值每次194.39 MiB，显式每次89.16 MiB（减少54.1\%）；进程工作集中位510.00变245.11 MiB（减少51.9\%）。普通选pressure、显式选stable，后者避免保存及验证多候选。仪器化墙时普通17.45--19.17秒、显式6.39--7.50秒，总64.94秒；\emph{不是普通吞吐基准}，不能把跟踪成本比写成全流程加速。证据\texttt{r74-bound-large-resources.json}保存四次源码、输入、输出哈希和全部尝试。它说明该限定工作负载的候选组合内存取舍，不改变普通默认质量或官方成绩。

\subsection{第七十五批：一般consumer下界与精确停止边界}
若一个原值有$u$个不同原消费者，设在最后消费前有$c$次该身份的COPY接力。从出生到最后消费的跨度至少包含$u+c$条指令；按时间排序的相邻身份接力（即使引用更早副本）及最后消费，每段距离不超过$W$，故跨度不超过$(c+1)W$。因此
\[
c\ \geq\ \max\!\left(0,\left\lceil\frac{u-W}{W-1}\right\rceil\right).
\]
对不同原身份求和，因为每条新增COPY仅属于一个身份。重复操作数在同一原指令中仅算一次消费者。它不依赖强制链、允许任意拓扑重排；当前实现仍保守筛选纯COPY-free输入、唯一且读写不相交的GET/SET寄存器。独立检查器从原JSON重新构建不同消费者集合，不读取调度器状态。240个至多10节点的纯小图与完整年龄状态oracle交叉核实，包括重复操作数及死节点；该下界不会证明别名、重算或额外GET模型的最优性。

R75新四家族、512/2048/8192/16384规模、四窗口、每格两种子，共128图，对普通auto、显式零容差、显式0.5\%停止三配置，总172.52秒。三版均104正确、16身份证书、8启发失败。零容差共同104图184706 COPY全相同，16例恰好达到独立全重排下界，其子集逐输入中位加速1.526；\emph{整体仅0.991}，p50由144.85变121.25 ms，p99由1128.82变1337.44 ms，不能把阳性子集包装成普遍收益，因此暂不默认。0.5\%版184719 COPY，增加13条，另列质量代价。104份严格输出换初值独立符号与三整数配置全部重放正确，14.21秒；原JSON数学下界另104份独立复核全部通过，13.14秒。输入清单SHA-256为\texttt{1be04bdb82506e872751a262a69d0833}\allowbreak\texttt{71a77289b34f73abe6dbeecaab9fc1d5}；证据\texttt{r75-consumer-bound}完整保存全部失败。

进一步尝试共同消费两个根值、各自仅由专用SET消费的$m$个计算。第一个和最后一个共同消费之间两根始终活跃；窗口最多容纳$W-2$个尚未存储的先前计算结果，所以该区间至少含
\[
L=\max(0,m-2)+\max(0,m-1-(W-2))
\]
条其他原指令。弱界用通用根年龄和范围$2(W-2)$。\emph{年龄坐标必须一致}：本文状态年龄按下一条待发射的位置计，\texttt{GET a}在0、\texttt{GET b}在1、\texttt{ADD(a,b)}在2，其ADD当前操作数距离为2、1（和3）；ADD发射后下一位置3的状态年龄为3、2（和5），不依赖额外链运算。若首共同计算在$p$、末共同计算在$q$，区间状态从$p+1$到$q$，恰含$q-p-1$个中间发射。加强版利用初态年龄至少2、3，末态在$q$操作数距离各不超过$W$且不同，至多$W-1,W$，端点差至多$2(W-3)$；根COPY一次至多减少年龄和$W-2$，给出
\[
C_{\rm pair}\geq\max\!\left(0,\left\lceil\frac{2L-2(W-3)}{W-2}\right\rceil\right).
\]
多个根组只能选\emph{不相交}根对替换逐身份贡献，不能对共享根重复累加。新增60配置、2--6共同消费者、至多14节点的小图完整oracle验证通过（只有五种规模拓扑，随机操作符、操作数顺序和SET顺序不能称60独立拓扑；专项显式扩大oracle节点上限，预算不足不算证明）。大规模新拓扑持出另批进行。它仍严格限定原节点身份，等值算术别名可能改变结果容量，因此不适用该新模型。

R76弱界的新32图包括随机1--4根组、512/2048/8192/16384规模、四窗口与两种子，三版全正确且83482 COPY相同，14次0.5\%停止，总181.66秒。相较普通auto，近似p50由964.32变172.23 ms，而逐输入中位只有1.004，p99由4440.28变4819.82 ms；不能只报分位数收益掩盖整体及长尾。32份严格输出另换初值独立符号/三整数重放全部通过，3.60秒。加强版专项8测试首次1.93秒、之后13.92秒的同机竞争差异也保留；各变体交替顺序，仍未完全控制外部负载。R77加强界另32图全部正确、68528 COPY一致，但\emph{没有一例}达到精确停止下界；整体逐输入中位1.002，总167.61秒。这说明更强必要条件仍不足以找到达到界限的实际发射。

\subsection{独立原节点身份证人验证}
原符号验证保持有序运算和计算多重集，但重复的等值原计算可能拥有相同SHA符号；语义等价不能单独证明输出没有复用另一个原身份。增加不信任的完整原节点排列证人，检查器仅从原JSON重新解析有序依赖，逐输出位置维护窗口内的\emph{整数原节点编号}：COPY继承源编号，每个非COPY行必须与下一原节点的操作符、寄存器和有序原编号参数完全一致，排列须是全部原位置恰一次。它不读取调度DAG、就绪集或内部年龄状态，也不依赖哈希相等。寄存器依赖独立按读取所处的前一SET epoch及SET写入顺序核实，允许无写屏障的读读交换；保守模式另要求全部GET/SET原序。原输入含COPY时无法仅凭非COPY行排列辨认原COPY，明确拒绝本项证书；本项数学质量界本来就排除此输入，普通语义校验仍支持它。

新增128不同家族/窗口小图、等值别名拒绝、读读许可和跨SET拒绝专项通过。

\texttt{milestone-v12}全部272成功输出重读身份证人/符号验证通过，4.79秒；历史R67、R68、R73、R75、R76、R77共344份严格结果全部通过，分别12.47、8.05、6.44、8.62、2.94、2.79秒，独立并行审计墙时不作为求解吞吐。\emph{两种目的分开}：符号及整数配置证明本地运算语义一致，整数身份重放证明质量界所需的原节点模型。只覆盖成功输出，不把未生成的模型证明、预算或启发失败算作输出验证。

R78将身份证人复核加入真实显式停止路径，并分别计入身份/数学检查时间。新48图包括共同根、脊柱及共享分支，512--16384规模、窗口4/32、两种子，旧新均48正确、936904 COPY相同；逐输入中位0.923，p50由243.13变206.22 ms而p99由2188.09变2285.99 ms。在线独立检查增加的成本和未控制CPU竞争均保留；不能拿离线检查成功掩盖线上成本。全部阳性停止现在同时具有原JSON身份重放和数学下界复核。

\subsection{严格窗口4叶分支的线性候选}
R79针对完整不相交双GET根组、每个ADD/MULW仅被一个专用SET消费、纯输入且唯一不相交读写寄存器的特殊形态，提出每组$m$计算的原序：前$\max(0,m-3)$个计算立即跟SET，最后至多三个计算集中后跟对应SET。窗口4时每个中间计算/SET之后需要两个根接力，尾部第三个计算结束两根均死亡；实际COPY为$\max(0,2m-6)$，等于加强端点下界。未使用的GET最后发射并保留。只窗口4、hazard模式、根对不相交、完整叶消费，其他结构正常回退；没有新增GET、删除运算、代数或等值别名。

32小中图三种原序提议中15图达到界限，窗口4的8图全部达到，其他窗口没有普遍收益。R79首次仅记录成本（6.05秒），另R79b固定同输入重跑补全输入哈希、严格输出及排列证人，1.10秒，二者分别保存而非覆盖。新160随机规模输入的线性候选均达到公式并通过独立语义/身份检查。候选只有实际COPY恰达原JSON复核的下界后才作\emph{零容差精确停止}；默认路径不会使用0.5\%近似。

\subsection{第八十、八十一批与最终冻结验证}
R80新32图、512/2048/8192/16384规模、四窗口、两种子，对冻结R78普通auto与加入窄结构候选的最终auto；两版32正确，8改善24相同，COPY83536变83506，无新失败或质量退化，总116.06秒。\emph{匹配W4的8图}均独立复核零容差最优，COPY54182变54152，逐输入中位10.277倍；\emph{整体32图中位仅1.058}，p50由724.02变310.34 ms、p99由4247.90变4059.63 ms，不能宣称一般10倍。其他窗口没有质量变化。32严格输出换初值三整数重放全部通过，4.78秒；独立整数身份重放14.35秒，后者并行CPU竞争较大，非求解吞吐。

R81独立种子持出64图、768/3072/6144/12288规模，包括匹配终端、一般共享根、随机树和稀疏共享。旧新均55正确、5容量证书、4启发失败；4改善51相同，COPY79664变79640，无新增失败，总71.01秒。匹配W4的4图均经在线原身份及数学复核达到零容差最优，COPY22200变22176，中位9.480倍；整体中位仅1.002，稀疏共享子集0.963，质量一致不能替代速度证据。55严格输出另三整数重放4.995秒、身份重放12.077秒全通过。仅已证明且可独立识别的W4完整形态进入默认，其他近似界、局部交换、重定位和分支实验继续显式。

最终\texttt{milestone-v14}在21拓扑、64/256/1024/4096规模、四窗口、每格一种子，共336图；旧新分别174/269输出，当前另47限定原身份模型容量证书、12预算耗尽、8启发失败，零错误输出。共同174图COPY20324变15213，40改善132相同2退化，新增95输出，逐输入中位6.389倍；总104.808秒。当前p50/p95/p99为32.978/256.838/287.081 ms；COPY两个等权分母均值为14.113\%（输出指令数）与42.621\%（原指令数），仅成功合成图，不能冒称官方达标。不同成功集合的总COPY不可直接相减。269严格输出重新读取后全部通过独立符号和三整数配置，10.398秒；原整数身份重放另6.743秒全通过。旧174缺原节点排列见证，\emph{仅符号/三整数通过、身份不可用}，不从新输出倒推旧见证。

v13评测适配曾对缺见证的旧版本直接索引\texttt{original\_order}，在符号通过后、数值检查前抛异常；旧结果被误计验证失败，该轮旧统计无效，原件完整保留，修复为明确身份不可用后新v14重跑336。所有线上质量证书仍要求完整身份见证，缺失即拒绝。v12/v13当前272输出与最终269的差异保留；其不匹配窄候选输入的核心组合并未改变，有限预算与并行机器竞争影响成功数，未以较高历史值替换最终成绩。

最终133回归全部通过，49.249秒；10000固定种子小图9998输出、2限定模型证明，7.618秒。两者在与v14完全相同的核心源码实际执行，最后只修复研究评测适配，复用已通过日志和报告，\emph{没有声称再执行一次}。来源、文件哈希、完整模块清单见\texttt{validation-provenance.json}。默认参数为auto、100 ms、30000状态、beam宽128、hazard；250 ms仅本轮对照预算。显式bound容差默认0.005，普通auto严格W4最优停止容差0。05:10冻结核心模块按相对路径排序JSON清单的SHA-256为\texttt{e27b01946a217034074345db2546867ce}\allowbreak\texttt{38ebbeb41a78258dbd888b14dc41d1f}；全部核心文件与v14启动时哈希一致。

R78在线身份成本整批51.75秒，48严格输出另三整数重放23.98秒、身份重放9.06秒全部通过；R77的32严格输出另三整数重放3.72秒通过。所有重放从保存的原输入、严格结果和不信任排列读入，不调用调度器，离线检查时间均与在线求解成本分开。

\subsection{测试硬件与单开发者工作量规划}
已测机器为Ryzen 9 8945HX（16核32线程）、31.3 GiB内存、Windows 11和CPython 3.12.14；当前算法、独立校验及单轮对照均单进程，未测GPU收益。8192节点的已测峰值工作集48.71--74.74 MiB只是四类r4输入，不是所有搜索的最大内存保证。规划最低硬件建议64位双核CPU、4 GiB内存、约5 GiB可用磁盘，可以做当前回归、小图和串行大图研究；建议8核CPU、16 GiB内存、10 GiB可用磁盘用于多参数/多block实验。两档均为工程建议，未在低配机器复测；精确状态预算需按内存限制设置。无GPU、开发板或商业求解器必需依赖。官方目标架构执行时间/功耗需取得对应模拟器或设备后另测，Python调度耗时不能替代目标ISA运行成本。

新增近窗口8192节点、bit\_width=5资源实测，插入121624 COPY，独立进程tracemalloc跟踪峰值114.49 MiB，进程峰值工作集315.44 MiB；含跟踪与独立检查墙钟12.73秒。该时间不能当作普通吞吐成绩；输出规模也须计入资源需求。证据\texttt{near-window-resources.json}保存源码/输入哈希。R74的16384原节点/24万以上COPY普通模式达到约510 MiB工作集，是截至该轮较高的实测例，\emph{不是所有搜索的内存上界}。05:10保存证据约3149971788字节（2.934 GiB），包括冻结源码、JSONL、严格输出、CSV/PDF及环境；磁盘建议须覆盖继续保存的证据，不能只按纯源码大小估计。

以下为一名熟悉编译器/图算法开发者、每天6--8小时的工程估算，不是已完成工时；官方说明到达后，需按真实规模重新估计。
\begin{table}[htbp]\centering\small
\begin{tabular}{p{0.49\linewidth}p{0.11\linewidth}p{0.30\linewidth}}\toprule
后续方向与验收 & 工作日 & 可并行部分\\\midrule
子图级重排/共享组处理：扩展已有局部原型，保留失败与退化 &3--5&拓扑生成与评测可并行\\
有界搜索状态表示与剪枝：独立oracle一致性与资源测量 &2--4&内存观测与验证审计可并行\\
官方格式/运算语义接入：严格导出与官方评分器对照 &1--3&与算法研究并行，依赖明确规范\\
CPU热点迁移C++（仅热点证据支持后）：双实现差分 &4--7&测试开发并行，内核开发顺序推进\\
Linux/低配复验与工程交付：一键运行和证据打包 &1--2&独立环境可并行\\\bottomrule
\end{tabular}\caption{后续工程估算，不代表上述方向都应立即启动。}\end{table}
首先推进3--5日的顺序优化及1--2日复验，形成下一可审查里程碑；约4--7工作日。各研究方向有依赖和重叠，不能直接把所有行相加作确定工期。多开发者可并行生成器、验证器、平台复验与资源测试；修改核心调度状态和组合策略应分阶段合并，避免共享状态或评测版本混淆。

\section{后续优化余量与研究判断}
首先应看同输入成功交集上的COPY变化和延迟，同时单列救回与退化；不同成功集合上的总COPY不能直接比大小。共享值和宽前沿可能超过窗口容量；GET延迟、子图完成顺序与提前接力可降低压力，但不能承诺所有DAG有解。较有希望的后续方向是搜索状态的年龄支配剪枝、窗口可行性引导的beam、独立子图分解与局部可验证修复；每种需要先明确是否保持完备性，再扩大固定输入对照。官方语义未定时，不应通过删除原节点、交换有序运算或额外读寄存器来制造看似优秀成绩。

测试成本可由输入数乘每个独立优化器调用的预算估计，但外部验证、预算检查粒度和Python运行时会增加墙钟。当前组合的250 ms是\emph{整个组合}的共享预算，不是每个候选都另获250 ms；例如384输入乘5个独立变体的工作预算为480秒。R30时十八家族、四规模、四窗口、每格两种子的一键批次为576输入，旧新两个变体名义工作预算288秒；R33增加第十九类后为608输入、304秒。零COPY提前结束会降低实际耗时，外部验证会增加时间。三计算模板的相同结构重复种子不增加有效结构样本量。上述批次单进程运行，但用户机器其他并行任务可能影响CPU竞争；跨时段p99不作严格因果比较。以后可并行独立basic block评测，须保留随机种子、输入哈希、相同预算及额外内存成本。

当前21家族、四规模、四窗口、每格两种子的一键批次为672输入，两个250 ms变体名义工作预算336秒；Samples=1为336输入、168秒。新别名家族是专门测模型边界的合成拓扑，不能代替普通应用分布。回归含比较穷尽oracle，运行成本不等于只调CLI若干次；新增语义模型与默认策略须单独保留源码快照及实验版本。

\begin{thebibliography}{9}
\bibitem{sugita} Shu Sugita, Toru Koizumi, Ryota Shioya, Hidetsugu Irie, Shuichi Sakai. A Sound and Complete Algorithm for Code Generation in Distance-Based ISA. CC 2023, pp.73--84. \url{https://doi.org/10.1145/3578360.3580263}. 作者机构记录：\url{https://pure.nitech.ac.jp/ja/publications/a-sound-and-complete-algorithm-for-code-generation-in-distance-ba/}. 作者组原文：\url{https://www.rsg.ci.i.u-tokyo.ac.jp/lab/en/assets/files/cc2023-sugita.pdf}.
\bibitem{sethi} Ravi Sethi, Jeffrey D. Ullman. The Generation of Optimal Code for Arithmetic Expressions. Journal of the ACM 17(4), 1970, pp.715--728. \url{https://doi.org/10.1145/321607.321620}.
\bibitem{llvm} LLVM Project. MachineScheduler.h API/source documentation. 访问日期2026-10-10. \url{https://llvm.org/doxygen/MachineScheduler_8h.html}.
\bibitem{hong} Jia-Wei Hong, H. T. Kung. I/O Complexity: The Red-Blue Pebble Game. STOC 1981, pp.326--333. 作者原文：\url{https://www.eecs.harvard.edu/~htk/publication/1981-stoc-hong-kung.pdf}.
\bibitem{unison} Roberto Casta\~neda Lozano, Mats Carlsson, Gabriel Hjort Blindell, Christian Schulte. Combinatorial Register Allocation and Instruction Scheduling. 作者预印本，2018，修订2019. \url{https://arxiv.org/abs/1804.02452}.
\bibitem{minor} Hans L. Bodlaender, Thomas Wolle, Arie M. C. A. Koster. Contraction and Treewidth Lower Bounds. Journal of Graph Algorithms and Applications, 10(1), 2006, pp.5--49. \url{https://doi.org/10.7155/jgaa.00117}. 期刊原文：\url{https://jgaa.info/index.php/jgaa/article/view/paper117}. 作者机构技术报告：\url{https://ics-archive.science.uu.nl/research/techreps/repo/CS-2004/2004-034.pdf}.
\end{thebibliography}
\end{document}
